\section{About the Topic}

\lead{The Baltic Sea has become an environment dedicated to a web of subsea cables (Figure 1), carrying vital telecommunications traffic between EU and NATO states. It also serves as a maritime link between Russia and Kaliningrad.\footnote{O’Riordan, K. (2025, April 11). Geopolitics of Baltic subsea infrastructure. Brussels Institute for Geopolitics. \url{https://big-europe.eu/publications/2025-04-11-geopolitics-of-baltic-subsea-infrastructure}}}

\guidefig[Submarine Cable Map \protect\footnotemark]{assets/figure-01.jpeg}[0.64]\footnotetext{Kivleniece, M. (2025). SEDE I: Cutting the cord—Undersea infrastructure security in the Baltic Sea. European Youth Parliament, The Hague 2025. \url{https://thehague2025.eyp.nl/to/sede-i/}}

These subsea cables, described as small and fiber-optic, span over 1,448,409.6 km of the seafloor from the Atlantic to the Pacific. They carry 99\% of the global internet data traffic, facilitating connection and communication across the globe. Despite their importance, these cables are fragile and highly prone to damage. \footnote{Himka, S. (2025, July 9). Baltic Sea undersea cable security. Henry M. Jackson School of International Studies, University of Washington. \url{https://jsis.washington.edu/news/baltic-sea-undersea-cable-security/}}

“The Baltic has become the decisive sea for NATO.” This was said by Commander

T. Gerhard Bidlingmaier, Federal German Navy, in September 1958. He also outlined that the region held great strategic importance back then for NATO’s northern and eastern flank. From this historical source, it is also obvious that the region has always been one of the checkpoints of tensions between NATO and the USSR/Russia. Throughout the Cold War, NATO has viewed control of the Baltic as important for limiting Russian/Soviet naval freedom of movement while keeping Allied access open. \footnote{Bidlingmaier, T. G. (1958, September). The strategic importance of the Baltic Sea. Proceedings, 84(9). \url{https://www.usni.org/magazines/proceedings/1958/september/strategic-importance-baltic-sea}}

In recent years, the world has been shaken by the growing number of incidents. Between November 2024 and January 2025 alone, seven cuts of the subsea cables have been registered. Suspicions were raised within the international community, with some speculating that the damage was not the result of sabotage. Due to the rapid growth in data volumes, subsea connections are becoming increasingly more important. Therefore, it can be expected that these critical connections will become targets of hostile acts as geopolitical conflicts intensify. \footnote{Muuga, Emilia \& Loik, Ramon \& Kaup, Georg-Henri \& Savimaa, Raul \& Koort, Erkki. (2025). Security Threats to the Undersea Connections Related Critical Infrastructure of the Baltic States: The Baltic Sea in the Focus of Hybrid Warfare. 10.15158/nv7t-kg46.} \footnote{Henry M. Jackson School of International Studies. (2025). Baltic Sea undersea cable security. University of Washington. \url{https://jsis.washington.edu/news/baltic-sea-undersea-cable-security/}}

Additionally, international bodies such as the EU and NATO make a careful distinction between hazards and accidents that occur in the region. There are around 200 accidental subsea cable faults every year. The majority of the faults stem from human activities, such as fishing. In \textbf{Figure 2} below, you can observe the source of these damages, according to the International Cable Protection Committee (ICPC). “Third Party” and “Unknown” labels include possible intentional sabotage.

\guidefig{assets/figure-02.jpeg}[0.43]

\textbf{Figure 2 - Distribution of causes for subsea cable damage (1959 - 2021)}\footnote{Pohoryles, D. A., Risks and protection of subsea cable networks, Publications Office of the European Union, 2025, p. 6, \url{https://data.europa.eu/doi/10.2760/0196933}} Last year, in March of 2025, the China Ship Scientific Research Center (CSSRC) and its affiliated State Key Laboratory of Deep-sea Manned Vehicles developed a ship able to cut cable lines at depths of up to 4,000 meters (13,123 feet). While China argues that the development of this equipment is purely for marine resource development, it still has caused further concern in the international community. \footnote{ibid} 13

ICPC has issued recommendations for accident prevention, including implementing monitoring systems and patrols. NATO has actually implemented flotilla patrols in the Baltic Sea. \footnote{ibid}

\subsection[Geography and natural resources]{Geography and natural resources\footnote{Murphy, E. L., \& Pearl, M., 2025, April, China’s underwater power play: The PRC’s new subsea cable-cutting ship spooks international security experts. Center for Strategic and International Studies. \url{https://www.csis.org/analysis/chinas-underwater-power-play-prcs-new-subsea-cable-cutting-ship-spoo} ks-international}}

The geographical aspect of the tension is crucial. The Baltic Sea is shallow and enclosed, accessible through three straits: the Great Belt, the Little Belt, and Øresund.

\textbf{Electricity:} the Baltic Sea region hosts numerous critical cross-border power cables connecting the Baltic states, Nordic countries, Poland, and Germany:

\begin{itemize}
  \item The NordBalt cable between Lithuania and Sweden,
  \item The Estlink 1 and Estlink 2 high-voltage direct current (HVDC) interconnectors between Finland and Estonia
  \item The SwePol cable between Sweden and Poland
  \item The Baltic Cable between Sweden and Germany
  \item Planned: Estlink 3 and Harmony Link interconnectors
\end{itemize}

\textbf{Gas:} the basin carries the Balticconnector pipeline between Finland and Estonia, and the Baltic Pipe from Norway through Denmark to Poland. The North Sea–Baltic data backbone runs through the C-Lion1 cable between Finland and Germany and the BCS East-West cable between Lithuania and Sweden.\footnote{Nordic Energy Research. (2026). Nordic energy security in two regional contexts: The Baltic and the Arctic. \url{https://pub.norden.org/nordicenergyresearch2026-02/section-4-nordic-energy-security-in-two-regional-contexts-the-baltic-and-the-arctic.html}}

\subsection{Vulnerability and danger}

The vulnerable nature of the cables, comprising extensive length and poor accessibility, makes them, unfortunately, easy and attractive targets for hybrid warfare. In late 2023, the Baltic Sea became a central theater of Russia’s undeclared hybrid war against NATO and its allies. It was the frequency, location, and methods of cable severance that pointed to a possible hybrid campaign, with Russia as the prime suspect. \footnote{International Centre for Defence and Security. (2025). The Baltic Sea in peace and war. \url{https://icds.ee/en/the-baltic-sea-in-peace-and-war/}}
